%%%%%%%%%%%%%%%%%%%%%%%%%%%%%%%%%%%%%%%%%
% Thesis 
% LaTeX Template
% Version 1.3 (21/12/12)
%
% This template has been downloaded from:
% http://www.latextemplates.com
%
% Original authors:
% Steven Gunn 
% http://users.ecs.soton.ac.uk/srg/softwaretools/document/templates/
% and
% Sunil Patel
% http://www.sunilpatel.co.uk/thesis-template/
%
% License:
% CC BY-NC-SA 3.0 (http://creativecommons.org/licenses/by-nc-sa/3.0/)
%
% Note:
% Make sure to edit document variables in the Thesis.cls file
%
%%%%%%%%%%%%%%%%%%%%%%%%%%%%%%%%%%%%%%%%%

%----------------------------------------------------------------------------------------
%	PACKAGES AND OTHER DOCUMENT CONFIGURATIONS
%----------------------------------------------------------------------------------------

\documentclass[11pt, a4paper, oneside]{Thesis} % Paper size, default font size and one-sided paper

\graphicspath{{./Pictures/}} % Specifies the directory where pictures are stored

\usepackage[square, comma, sort&compress]{natbib} % Use the natbib reference package - read up on this to edit the reference style; if you want text (e.g. Smith et al., 2012) for the in-text references (instead of numbers), remove 'numbers'
\usepackage[utf8]{inputenc}
\usepackage[linesnumbered,ruled]{algorithm2e}
\usepackage{amsmath}
\hypersetup{urlcolor=black, colorlinks=true} % Colors hyperlinks in blue - change to black if annoying
\title{\ttitle} % Defines the thesis title - don't touch this

\begin{document}

\frontmatter % Use roman page numbering style (i, ii, iii, iv...) for the pre-content pages

\setstretch{1.3} % Line spacing of 1.3

% Define the page headers using the FancyHdr package and set up for one-sided printing
\fancyhead{} % Clears all page headers and footers
\rhead{\thepage} % Sets the right side header to show the page number
\lhead{} % Clears the left side page header

\pagestyle{fancy} % Finally, use the "fancy" page style to implement the FancyHdr headers

\newcommand{\HRule}{\rule{\linewidth}{0.5mm}} % New command to make the lines in the title page

% PDF meta-data
\hypersetup{pdftitle={\ttitle}}
\hypersetup{pdfsubject=\subjectname}
\hypersetup{pdfauthor=\authornames}
\hypersetup{pdfkeywords=\keywordnames}

%----------------------------------------------------------------------------------------
%	TITLE PAGE
%----------------------------------------------------------------------------------------

\begin{titlepage}
\begin{center}

\textsc{\LARGE \univname}\\[0.1cm]
\textsc{Faculteit der Natuurwetenschappen, Wiskunde en Informatica} \\[1cm] % University name
\textsc{\Large MSc Thesis}\\[0.5cm] % Thesis type

\HRule \\[0.4cm] % Horizontal line
{\huge \bfseries \ttitle}\\[0cm] % Thesis title
\HRule \\[1.5cm] % Horizontal line

\begin{minipage}{0.4\textwidth}
\begin{flushleft} \large
\emph{Author:}\\
\authornames % Author name - remove the \href bracket to remove the link
\end{flushleft}
\end{minipage}
\begin{minipage}{0.4\textwidth}
\begin{flushright} \large
\emph{Supervisor:} \\
\supname % Supervisor name - remove the \href bracket to remove the link  
\end{flushright}
\end{minipage}\\[3cm]
 
\large \textit{A thesis submitted in fulfillment of the requirements\\ for the degree of \degreename}\\[0.3cm] % University requirement text
\textit{in the}\\[0.4cm]
\deptname\\[2cm] % Research group name and department name

{\large \today}\\[2cm] % Date
\includegraphics[height=3cm]{uva-logo} % University/department logo - uncomment to place it

\vfill
\end{center}

\end{titlepage}

%----------------------------------------------------------------------------------------
%	DECLARATION PAGE
%	Your institution may give you a different text to place here
%----------------------------------------------------------------------------------------

%\Declaration{

%\addtocontents{toc}{\vspace{1em}} % Add a gap in the Contents, for aesthetics

%I, \authornames, declare that this thesis titled, '\ttitle' and the work presented in it are my own. I confirm that:

%\begin{itemize} 
%\item[\tiny{$\blacksquare$}] This work was done wholly or mainly while in candidature for a research degree at this University.
%\item[\tiny{$\blacksquare$}] Where any part of this thesis has previously been submitted for a degree or any other qualification at this University or any other institution, this has been clearly stated.
%\item[\tiny{$\blacksquare$}] Where I have consulted the published work of others, this is always clearly attributed.
%\item[\tiny{$\blacksquare$}] Where I have quoted from the work of others, the source is always given. With the exception of such quotations, this thesis is entirely my own work.
%\item[\tiny{$\blacksquare$}] I have acknowledged all main sources of help.
%\item[\tiny{$\blacksquare$}] Where the thesis is based on work done by myself jointly with others, I have made clear exactly what was done by others and what I have contributed myself.\\
%\end{itemize}
 
%Signed:\\
%\rule[1em]{25em}{0.5pt} % This prints a line for the signature
 
%Date:\\
%\rule[1em]{25em}{0.5pt} % This prints a line to write the date
%}

%\clearpage % Start a new page

%----------------------------------------------------------------------------------------
%	QUOTATION PAGE
%----------------------------------------------------------------------------------------

\pagestyle{empty} % No headers or footers for the following pages

\null\vfill % Add some space to move the quote down the page a bit

\textit{``Thanks to my solid academic training, today I can write hundreds of words on virtually any topic without possessing a shred of information, which is how I got a good job in journalism."}

\begin{flushright}
Dave Barry
\end{flushright}

\vfill\vfill\vfill\vfill\vfill\vfill\null % Add some space at the bottom to position the quote just right

\clearpage % Start a new page

%----------------------------------------------------------------------------------------
%	ABSTRACT PAGE
%----------------------------------------------------------------------------------------

\addtotoc{Abstract} % Add the "Abstract" page entry to the Contents

\abstract{\addtocontents{toc}{\vspace{1em}} % Add a gap in the Contents, for aesthetics

%From the project outline
A basic fact about human-human dialogue is that there is often more than one way of talking about any domain. For example, instead of saying \textit{``Turn right after 200 meters"}, a route giver may say \textit{``Turn right at Barclays"} given that \textit{``Barklays"} is an enstabilished referring expression. Human speakers efficiently align terminology and associated ontology in talking about any specific domain (in this example we can say that ontologically there is only one entity -- a location -- but there are different linguistic terms we can use to refer to it). In contrast, current dialogue systems typically have a static ontology and a static vocabulary, which is used in both generation and interpretation, requiring users to formulate their utterances using the terminology known by the system. The goal of this project is to work towards a system that adapts its own linguistic resources (including possible ontology, vocabulary and grammar) to the interlocutor. The system should be able to learn new concepts by assigning  new meanings to known words, as well as new words to talk about concepts known by the system in the domain.
}

\clearpage % Start a new page

%----------------------------------------------------------------------------------------
%	ACKNOWLEDGEMENTS
%----------------------------------------------------------------------------------------

\setstretch{1.3} % Reset the line-spacing to 1.3 for body text (if it has changed)

\acknowledgements{\addtocontents{toc}{\vspace{1em}} % Add a gap in the Contents, for aesthetics

The acknowledgements and the people to thank go here\ldots

% Other people to acknowledge
% - The guy who created this template: http://www.latextemplates.com/template/masters-doctoral-thesis
}
\clearpage % Start a new page

%----------------------------------------------------------------------------------------
%	LIST OF CONTENTS/FIGURES/TABLES PAGES
%----------------------------------------------------------------------------------------

\pagestyle{fancy} % The page style headers have been "empty" all this time, now use the "fancy" headers as defined before to bring them back

\lhead{\emph{Contents}} % Set the left side page header to "Contents"
\tableofcontents % Write out the Table of Contents

\lhead{\emph{List of Figures}} % Set the left side page header to "List of Figures"
\listoffigures % Write out the List of Figures

\lhead{\emph{List of Tables}} % Set the left side page header to "List of Tables"
\listoftables % Write out the List of Tables

%----------------------------------------------------------------------------------------
%	ABBREVIATIONS
%----------------------------------------------------------------------------------------

\clearpage % Start a new page

\setstretch{1.5} % Set the line spacing to 1.5, this makes the following tables easier to read

\lhead{\emph{Abbreviations}} % Set the left side page header to "Abbreviations"
\listofsymbols{ll} % Include a list of Abbreviations (a table of two columns)
{
\textbf{AI} & \textbf{A}rtificial \textbf{I}ntelligence \\
\textbf{ASR} & \textbf{A}utomatic \textbf{S}peech \textbf{R}ecognition \\
\textbf{CA} & \textbf{C}onversational \textbf{A}gent \\
\textbf{DM} & \textbf{D}ialogue \textbf{M}anager \\
\textbf{ISU} & \textbf{I}nformation \textbf{S}tate \textbf{U}pdate \\
\textbf{MT} & \textbf{M}achine \textbf{T}ranslation \\
\textbf{NLP} & \textbf{N}atural \textbf{L}anguage \textbf{P}rocessing \\
\textbf{SLU} & \textbf{S}poken \textbf{L}anguage \textbf{U}nderstanding \\
\textbf{TTS} & \textbf{T}ext \textbf{T}o \textbf{S}peech \\
%\textbf{Acronym} & \textbf{W}hat (it) \textbf{S}tands \textbf{F}or \\
}

%----------------------------------------------------------------------------------------
%	SYMBOLS
%----------------------------------------------------------------------------------------

\clearpage % Start a new page

\lhead{\emph{Symbols}} % Set the left side page header to "Symbols"

\listofnomenclature{lll} % Include a list of Symbols (a three column table)
{
$c$ & chunk of text \\
$m$ & meaning \\
$s$ & sentence \\
% Symbol & Name & Unit \\

& & \\ % Gap to separate the Roman symbols from the Greek

$\sigma$ & chunk similarity \\
$\tau$ & word similarity \\
% Symbol & Name & Unit \\
}

%----------------------------------------------------------------------------------------
%	DEDICATION
%----------------------------------------------------------------------------------------

%\setstretch{1.3} % Return the line spacing back to 1.3

%\pagestyle{empty} % Page style needs to be empty for this page

%\dedicatory{For/Dedicated to/To my\ldots} % Dedication text

%\addtocontents{toc}{\vspace{2em}} % Add a gap in the Contents, for aesthetics

%----------------------------------------------------------------------------------------
%	THESIS CONTENT - CHAPTERS
%----------------------------------------------------------------------------------------

\mainmatter % Begin numeric (1,2,3...) page numbering

\pagestyle{fancy} % Return the page headers back to the "fancy" style

% Include the chapters of the thesis as separate files from the Chapters folder
% Uncomment the lines as you write the chapters

% Chapter Template

\chapter*{Introduction} % Main chapter title

\label{ch:Introduction} % Change X to a consecutive number; for referencing this chapter elsewhere, use \ref{ChapterX}

\lhead{\emph{Introduction}} % Change X to a consecutive number; this is for the header on each page - perhaps a shortened title

%----------------------------------------------------------------------------------------
%	INTRO TEXT
%----------------------------------------------------------------------------------------


The \textbf{user interface}, or human-computer interface (HCI) is the component of a computer system that provides a space of interaction between the human user and the resources offered by the machine; such a space defines a bridge language which human intentions can be translated into, to be converted into computational procedures for the machine; vice versa, the result of the computation is then presented to the user in the same language, which he or she is assumed to understand.

% cite Tanenbaum
In the early days of computing, the so-called \textbf{batch interfaces} were non-interactive: the user/programmer was supposed to feed the machine with a software, punched on cards using the \textit{machine's assembly language} directly, and retrieve the result of the computation printed on paper. The third and fourth generations of computing brought \textbf{text-based} interfaces for operating systems (UNIX, DOS, CP/M), that were later took over by \textbf{Graphical User Interfaces} (GUI), moving the human-machine interaction on a new, visual language made of windows, icons and buttons. Recently, touch screen and camera devices allowed for the implementation of even more natural means of interaction based on \textbf{gestures} and physical actions.

From this brief spot on computing history, and in a way from common sense, we can draw the rather trivial, yet crucial, conclusion that the trend in user interfaces development is to \textbf{close the gap} between humans and machines by moving the needle of the interface languages from a machine-centered space towards the human language itself. To this respect, the studies on Natural Language Processing (NLP) assume a dramatically central role, as dramatically central is natural language in the interaction of humans with each other.

%----------------------------------------------------------------------------------------
%	DIALOGUE SYSTEMS
%----------------------------------------------------------------------------------------

\section*{Dialogue Systems}

A spoken dialogue system, or conversational agent (CA), allow humans and machines to interact through an intermediate language which is as close as possible to the \textbf{human language}, and through conversational episodes that implement as close as possible the human dialogue modalities.

Research in dialogue systems has been carried on since the \textbf{early days} of Artificial Intelligence. A milestone in the early work on this field is ELIZA \cite{Weizenbaum:1966:ECP:365153.365168}, which provides the user with a basic human-like interaction based on pattern matching; another example is the SHRDLU system
% Cite SHRDLU: http://dspace.mit.edu/handle/1721.1/7095?show=full
, which interfaces the user with a simple spatial domain being able to ambiguous or implicit references to the entities in it.

According to \cite{Jokinen2009}, modern dialogue systems can be divided in \textbf{two main types}: task-oriented and nontask-oriented. Intuitively, systems in the first category are meant to deal with a specific task such as making a hotel booking, or booking a plane ticket; an example in this category is the MIT Mercury system, a vocal interface to a flight database \cite{Seneff:2000:DMM:1605285.1605288}. On the other hand, nontask-oriented systems are meant to engage in conversations without a specific purpose to fulfill, but the one of delivering a realistic simulation; ELIZA itself is an example of nontask-oriented dialogue system.

Task-oriented systems can be very simple, as simple and well-formalized the task is;  many applications, such as travel service or call routing, can be successfully solved by \textbf{slot-based} systems: each step of the conversation requires some pieces of information, modeled as slots, to be filled in by the user (departure city, arrival city, date, and so on); given the slots to be filled, the dialogue task can be solved with a formal grammar of interaction. As the complexity increases, more phenomena of human interaction have to be modeled, such as turn-taking, multimodality or grounding
%\ignore{CITE RAQUEL'S CHAPTER}
, as well as semantic structures such as quantification and negation; slot-based systems are not sufficient to model these scenarios \cite{Gabsdil03clarificationin}, that require more advanced frameworks such a the Information State Update (ISU) one.

%----------------------------------------------------------------------------------------
%	LEARNING TO TALK
%----------------------------------------------------------------------------------------

\section*{Learning to talk}

One of the constituent features of humans' ability to speak is that such an ability is \textbf{not innate}, but is rather learned through interactions.

Ever since the power of computers grew enough to allow for intensive statistical analysis of significant amounts of data, \textbf{Machine Learning} approaches to Artificial Intelligence tasks got more and more prominent in the scene, often outperforming static methods (i.e. where the solution procedure for a task is explicitly coded by the programmer).
% Citation needed...
One of the clearest examples is the field of NLP itself, where the most important tasks, like parsing or machine translation, are dominated by Machine Learning methods based on corpora,
% Citation needed
meaning that, for instance, a Machine Translation system will first be trained on a corpus of aligned sentences. Statistical structures will be extracted from this corpus, like the most likely word-by-word alignment, and later be used to process new examples.

However, the task of learning dialogue interactions brings some \textbf{peculiar} challenges. First of all, humans learning to talk do not go through two separate phases of learning and processing, but rather improve their abilities episode by episode; as \cite{2095408} point out, this incremental learning structure is nowadays not implemented in state-of-the-art systems. Furthermore, ...
\begin{itemize}
	\item Learning to produce or understand new surface forms, or \textbf{realizations}, for a given meaning -- eg. the sentence ``Bill eats an apple" for the action of Bill eating an apple.
	\item Learning to produce new \textbf{meanings} from the existing ones and their respective realizations -- eg. the concept of motor home, sharing the features of a house and a car.
	\item Learning a \textbf{grammar of conversation}, to place the correct utterance at each step of a conversational episode. -- eg. an appropriate answer to the utterance ``My name is Bill" is ``Nice to meet you", whereas ``I like cookies" would not sound as much appropriate.
\end{itemize}
Lastly, we can point out that the definitions of actions, meaning and episodes are somehow arbitrary and not sharply bounded, as it can be argued for recursive and compositional structures at any level of their interpretation.

%----------------------------------------------------------------------------------------
%	THIS THESIS PROJECT
%----------------------------------------------------------------------------------------

\section*{This thesis project} \label{ch:intro:project}

The aim of this thesis project is to design and implement language learning capabilities for an existing dialogue system, focusing on the \textbf{realization level}. That is, given a fixed list of meanings, the system should be able to classify every sentence into its correct meaning.

The \textbf{domain} I am considering is the one of a music player application, being able to get natural language input from the user and translate it into an appropriate corresponding behaviour. For example, when the input is ``Play Pictures at an exhibition", the system should start playing the famous suite by Modest Mussorgsky.

In this domain, each \textbf{meaning} is an action that a player is performed, and is defined by a set of representative sentences, being its surface forms. For instance, the action of increasing the volume level can be defined by the following set of sentences:
\begin{itemize}
	\item \textit{Increase the volume}
	\item \textit{Increase the volume level}
	\item \textit{Raise the volume}
	\item \textit{Increase the volume please}
\end{itemize}

Note that, even though one would be likely to think of such meanings in compositional terms, such compositionality have not been explored for the sake of simplicity.

The \textbf{task} for the application is, given an arbitrary input sentence and a context (the point of the conversational episode being realized), to perform the correct action, that is, associate that sentence with its correct meaning. Furthermore, the system should be able to \textbf{learn} new realizations for each meaning, as unknown sentences are given in input and processed.

Note that such processing might be more or less \textbf{semantically intensive}. As an example, it can be argued that, given the above definition of the action to increase the volume, matching \textit{``raise the volume please"} is an easier task than classifying \textit{``Turn up the volume"}.

Also, an unknown input sentence should be given a \textbf{confidence score} for each candidate meaning it is associated to; the system should be able to narrow possible needs for clarification down to single sentence components, eventually asking the user for disambiguation as specifically as possible.

This document is structured as follows. Chapter \ldots
% Chapter Template

\chapter{Architecture and related work} % Main chapter title

\label{ch:arch_rw} % Change X to a consecutive number; for referencing this chapter elsewhere, use \ref{ChapterX}

\lhead{Chapter \ref{ch:arch_rw}. \emph{Architecture and related work}} % Change X to a consecutive number; this is for the header on each page - perhaps a shortened title

%----------------------------------------------------------------------------------------
%	SECTION 1
%----------------------------------------------------------------------------------------

\section{Architecture}

\ldots 

\chapter{The M2 matching algorithm} % Main chapter title

\label{ch:M2} % Change X to a consecutive number; for referencing this chapter elsewhere, use \ref{ChapterX}

\lhead{Chapter \ref{ch:M2}. \emph{The M2 matching algorithm}} % Change X to a consecutive number; this is for the header on each page - perhaps a shortened title

Computing a \textbf{similarity score} between two natural language \textbf{sentences} is not a trivial problem. Many features can be taken into account to solve it, such as the syntactic structure of the sentence, or the meaning of the single words, or their order; all of this features are informative respect to the meaning of the sentence, and it is reasonable to assume they are all considered by humans solving the same task. Here I describe a recursive, dynamic programming \textbf{algorithm} for sentence comparison aimed to exploit word-based information, as well as the syntactic one, expressed in the form of unsupervised learned binary trees.

\section{Introduction}
In the context of sentences \textbf{classification}, where the task is to label an unknown input sentence with the correct meaning, and each meaning is defined by a set of representative natural language sentences (as it was defined in \ref{ch:intro:project}), one crucial point is to measure the similarity of the input sentence with each of the sentences defining each meaning.

The problem of scoring the similarity between two sentences is not new in the \textbf{literature}, and a number of different approaches already exist to tackle it. Achananuparp et al. in \cite{achananuparp08} classify these approaches in \emph{word overlap} measures (based on the number of words shared by the two sentences), \emph{TF-IDF} measures (based on term frequency-inverse document frequency) and \emph{linguistic} measures (based on semantic relations between words and their syntactic composition).

Another source of inspiration for this work is represented by \textbf{statistic}, corpus-based methods in Computational Linguistics; a significant example comes from The IBM models for Statistical \textbf{Machine Translation} \cite{ibmmt}, that first introduced the idea of feeding statistically intensive \textbf{Machine Learning} algorithms with big data from corpora, which nowadays is the dominant paradigm in MT; insightful is also the work on Data Oriented Parsing, and particularly the U-DOP model for \textbf{Unsupervised Language Learning} \cite{udop}, which core idea is to initially assume all the possible syntax trees for a set of sentences as equally possible, and then use all the possible sub-trees of them to compute the most probable parse trees, letting the structure of the language emerge from the data.

% A bit more on literature, relation to the work would be nice..

The algorithm presented in this paper aims to provide a similarity measure for two sentences $s_1$ and $s_2$ featuring many of the ideas that have proven to be successful in recent developments of Computational Linguistics. The following are some of the key insights of the algorithm:

\begin{itemize}
\item \textbf{Every possible} sub-sentence of $s_1$ is compared with every possible sub-sentence of $s_2$.
\item A score between two chunks of text is build \textbf{recursively} on the best matches of smaller chunks, and always brought down to scores between single words.
\item \textbf{Dynamic programming} is used to increase the efficiency of the algorithm, preventing it from computing the same result two times.
\item Other than the similarity score, the algorithm outputs the most probable chunking of each sentence in the form of an unlabeled binary \textbf{parse tree}.
\item Other than the similarity score and the parse trees, the algorithm outputs the most likely \textbf{alignment} between the chunks of the two sentences.
\item Every score is expressed as a linear combination of \textbf{features}, designed to be independent and extensible.
\item The partial computations of the algorithm can be stored and further used to implement a learning model.
\end{itemize}

%Document structure
This paper is \textbf{structured} as follows. Section \ref{algo} describes the algorithm theoretically. %and in a better world a complexity analysis
%Section \ref{example} walks through a toy example.
Section \ref{learning} describes possible extensions and improvements. Section \ref{conclusions} draws the conclusion of the paper.

\section{The algorithm} \label{algo}
\begin{algorithm}
  \SetKwData{Left}{left}
  \SetKwData{Up}{up}
  \SetKwFunction{WWScore}{$\tau$}
  \SetKwFunction{CCScore}{$\tilde{\sigma}$}
  \SetKwFunction{Length}{Length}
  \SetKwFunction{Split}{Split}
  \SetKwFunction{Set}{Set}
  \SetKwFunction{Append}{Append}
  \SetKwFunction{Max}{Max}
  \SetKwInOut{Input}{input}
  \SetKwInOut{Output}{output}
  \SetKw{In}{in}

\Indm  
  \Input{$c_1,c_2$ chunks of text; $T$ partial results table}
  \Output{$\sigma(c_1,c_2)$ Similarity score between $c_1$ e $c_2$}
\Indp
  \BlankLine
  \uIf{c1,c2 are words}{ \label{main:bc1}
	  \Return{\WWScore{$c_1,c_2$} \label{main:bc2} }
  }
  \BlankLine
  $S\leftarrow[\ ]$\; \label{main:cinit}
  \For{$i\leftarrow 1$ \KwTo \Length{$c_1$}}{ \label{main:loopi}
    $C_1 \leftarrow $ \Split{$c_1$,i}\; \label{main:spliti}
		\For{$j\leftarrow 1$ \KwTo \Length{$c_2$}}{ \label{main:loopj}
		\lIf{i=\Length{$c_1$} and j=\Length{$c_2$}}{break} \label{main:break}
		\BlankLine
		$C_2 \leftarrow $ \Split{$c_2$,j}\; \label{main:splitj}
		\BlankLine
		\ForEach{$sc_1$ \In $C_1$}{ \label{main:fe1}
			\ForEach{$sc_2$ \In $C_2$}{ \label{main:fe2}
				\uIf{$T[sc_1,sc_2]=\emptyset$}{$T[sc_1,sc_2] \leftarrow \sigma(sc_1,sc_2,T)$ \label{main:tabass} \label{main:tupdate}}
			}
		}
		\BlankLine
		\Append{$S$,\CCScore{$C_1,C_2,T$}} \label{main:cappend}
 }
	}
	\BlankLine
	\Return \Max{$R$} \label{main:return}
\caption{The main algorithm\label{main}}
\end{algorithm}

Algorithm \ref{main} contains the main loop of the sentence scoring algorithm. Before entering the details of it, it is necessary to say that:
\begin{itemize}
\item The input $c_1,c_2$ can be any \textbf{chunks} of text because, while the algorithm is called with two sentences in input, it will recursively call itself on every possible chunk extracted from the two sentences.
\item The input $T$ (T for Table) is a data structure holding the intermediate results of the algorithm. As we can infer from line \ref{main:tabass} of Algorithm \ref{main}, $T[c_1,c_2]=\sigma(c_1,c_2)$
\item $\tau(c_1,c_2)$ returns a similarity score between two words. This score is a linear combination of independent features. At the moment, the following features are taken into account:
\begin{itemize}
\item Equality (boolean) -- 1 if $c_1=c_2$, 0 otherwise
\item Edit distance -- $1-d$, where $d$ is the Levelshtein distance \cite{levelshtein-66-binary} between $c_1$ and $c_2$, as it is implemented in NLTK \cite{Loper:2002:NNL:1118108.1118117}
\item Difference in position -- $(1/|p_{c_1}-p_{c_2}|)^\alpha$, where $p_{c_i}$ is the position of $c_i$ in the sentence, and $\alpha$ is a free parameter.
\item Wordnet \cite{Miller:1995:WLD:219717.219748} Path Length similarity, as it is implemented in NLTK
\end{itemize}
The weights of the single features, now static, are meant to be trained through Machine Learning
\item \texttt{Length($c$)} returns the number of words contained in the chunk $c$
\item \texttt{Split($c,i$)} returns a list of two chunks, one containing the words of $c$ from the first to the $i$th, and the other containing the words of $c$ form the $(i+1)$th to the last.

If the $i$th word is the last word of $c$, it returns a list containing the only element $c$.
\item \texttt{Append($L,x$)} is the standard \emph{append} operation for lists, adding the element $x$ in the last position of the list $L$.
\item $\hat{\sigma}(C_1,C_2,T)$ returns the score of $c_1,c_2$ \textbf{given} a particular 2-split of the two chunks. Such a score is computed combining the ones of the smaller constituents of $c_1,c_2$, which are required to be already present in the table T.

The score is again the linear combination of independent features, leaving its implementation open and extensible.

At the moment the only feature being employed is the average of the scores of the best alignments of the sub-chunks in $C_1$ with the ones in $C_2$.

For example, given $C_1=[\text{"increase","the volume"}]$ and $C_2=[\text{"raise"},\\ \text{"the volume"}]$, "increase" is likely to achieve the best score with "raise" (eg. $T[\text{"increase","raise"}]=0.8$), as well as "the volume" will be aligned with "the volume" ($T[\text{"the volume","the volume"}]=1$)). The score of "increase the volume" and "raise the volume", in this particular splitting, will thus be the average of the two, 0.9.
\end{itemize}
After these premises it is relatively easy to walk through the pseudo-code of Algorithm \ref{main}. Lines \ref{main:bc1} and \ref{main:bc2} handle the base case, that is, when the algorithm is run on \textbf{two words}; in this case the result of the Word Similarity function, which is described above, is returned.

At line \ref{main:cinit} the list of \textbf{candidate results} is initialized as an empty list. This is because more than one score will be computed for $c_1$ and $c_2$, the maximum of which will be returned as a result.

Lines \ref{main:loopi} and \ref{main:loopj} define a nested loop structure, which purpose is to update the two indexes $i$ (from the first to the last word of $c_1$) and $j$ (from the first to the last word of $c_2$) to cover \textbf{every possible combination} of word positions in the two input strings. These indexes are used to \textbf{split} the input chunks in smaller parts (lines \ref{main:spliti} and \ref{main:splitj}). As an example, if the input is given by $c_1=$"increase the volume", $c_2=$"raise the volume", the outer loop will produce the three possible values of $C_1$: ["increase","the volume"], ["increase the","volume"] and ["increase the volume"]; each of these values will be combined, in the inner loop, with the three possible values of $C_2$: ["raise","the volume"], ["raise the","volume"] and ["raise the volume"].

Note that, while it is reasonable to consider the entire $c_1$ and $c_2$ in the comparisons (eg. if $c_1=$"quit",$c_2=$"quit please", we'd want the whole $c_1$ to be associated with the "quit" sub-part of $c_2$), it is necessary to prevent the two entire $c_1$ and $c_2$ to be associated at the same time, thus producing an \textbf{infinite loop}; this is done in line \ref{main:break}.

The next part of the inner loop, from line \ref{main:fe1} to line \ref{main:fe2} combines every sub-chunk of $c_1$ ($sc_1 \in C_1$) with every sub-chunk of $c_2$ ($sc_2 \in C_2$); note that $C_1$ and $C_2$ contain either one or two elements. Every combination that is not present in the table is scored with a \textbf{recursive} call to the algorithm itself, and the resulting score is saved in the table (line \ref{main:tupdate}).

The last operation done in the loop, at line \ref{main:cappend}, is to combine the sub-scores that have just been scored in the table to compute the final score of this particular subdivision of $c_1$ and $c_2$. This is a \textbf{candidate score} for $c_1$ and $c_2$, and thus is appended to the list of candidate scores. As we already said, the definition of this combination is open; however it is most reasonable for the $\hat{\sigma}$ function to find the \textbf{best alignment} between the input sub-chunks, and provide a score based on it. It is clear that the definition of $\hat{\sigma}$ is crucial for getting sensible results from the algorithm. Furthermore, $\hat{\sigma}$ is the point where the algorithm can be extended with Machine Learning capabilities, that will be discussed in Section \ref{learning}.

At line \ref{main:return} the function returns the \textbf{maximum of the candidate} score, that is, the split that produced the best score (eg. for "increase the volume" vs "raise the volume", the score of "increase $\vert$ the volume" vs "raise $\vert$ the volume" is likely to produce a better score than "increase $\vert$ the volume" vs. "raise the $\vert$ volume"). Note that this decision determines the branching of one level of the recursion tree; once the algorithm is done computing the score between two sentences, a traceback of the selected splits can be interpreted as a parse tree for the two input sentences.

%\section{An example} \label{example}

\section{Machine Learning Possibilities} \label{learning}
The basic algorithm described in section \ref{algo} only deals with \textbf{static features} to measure similarities: the score of two words is derived from equality, edit distance, Wordnet path length and so on; the score of chunks is the average of the scores of the best alignments of their components.

In this section I address the question whether it is possible, and how, for the algorithm to \textbf{learn} from the sentences it is parsing. In other words, given a sentence of which I know the meaning, how can I use this information to improve the way I process new sentences?

Here I consider \textbf{three features} based on Statistical Language Processing methods to address this problem: chunk likelihood, class-conditional chunk weight and alignments likelihood.

\subsection{Chunk likelihood}
% Relative frequency across all sentences
% Used to select good chunks
Let's consider the following pair of sentences:
\begin{quote}
Pump up the volume\\
Turn up the volume
\end{quote}

As an English-speaking human being I immediately associate the phrase "pump up" with "turn up" and "the volume" with its analogue. One of the reasons I make this association rather than, for instance, associating "pump" with "turn" and "up the volume" with its analogue, is that I have good experience of the phrase "the volume" being used in different contexts, while fewer times I encounter "up the volume".

Thus a feature that can be useful to \textbf{enforce good chunking} of the input sentences is the likelihood of a chunk of text (where a chunk is defined as either a word or a combination of chunks). This feature can be modeled as the relative count of each chunk in the whole pool of the ones that have been recorded:
$$
l(c)=\frac{\text{\# of }c}{\text{\# of total chunks}}
$$

Note that this feature is completely unsupervised, in that it does not require a Meaning label to be computed.
\subsection{Class-conditional chunk weight}
% frequency in class / frequency overall
% Used to weight more chunks that are representative for one class
Let's consider the following three sentences:
\begin{quote}
Increase the volume\\
I would appreciate if you could increase a bit the the volume\\
Decrease the volume
\end{quote}

With a reasonable approximation, the first two sentences have the same meaning, while the last says the opposite, although the first and the last sentence are much more similar from a syntactic point of view. Nevertheless, no English speaker would bother wondering about hidden meanings behind phrases like "I would appreciate", or "if you could": everyone could tell that the second sentence is no more than the first one, with some more formal dressing.

Especially, when I listen to the second sentence, I can immediately locate "increase" and "the volume" as the most informative phrases, while filtering out the rest as less relevant. One of the reasons I am able to do this is that I have memory of "I would appreciate if you could" used in sentences conveying lots of different meanings, while "increase" is likely to appear only in situations where something is increased.

A feature that could be useful to spot the \textbf{informativeness} of chunks can be the conditional probability of the chunk, given a certain meaning:
$$
l(c|m)=\frac{\text{\# of }c\text{ in }m}{\text{\# of }c}
$$

This feature can be used to weight the score of different chunks in the sentence, so for more informative chunks to influence more of the total score mass.
%l(c|m,M)?
\subsection{Alignment likelihood}
% Recursive score of alignments
% Used to align chunks correctly
Let's consider the following pair of sentences:
\begin{quote}
Skip this track \\
Jump over this song
\end{quote}
Again, English speakers can easily associate the meaning of "skip" with the one of "jump over", and the meaning of "track" with the one of "song".

One device that the algorithm uses to solve this problem is the \textbf{Wordnet} Path Lenght measure for word comparison. However, it has been noticed that this measure has its limits. In particular
\begin{itemize}
\item Not all the words are included in Wordnet
\item Being a static measure, it cannot be adapted to the domain
\item It has difficulties spotting antonyms
\end{itemize}
The alignment likelihood measure can be used to \textbf{enforce correct alignments} of different chunks.
$$
l(c_2|c_1)=\frac{\text{\# of }c_1-c_2\text{ alignments }}{\text{\# of }c_1\text{ alignments}}
$$
Note that, since no alignment labels are provided in the dataset (each sentence is labeled with a meaning, but its syntactic structure is hidden), this measure have to be trained in an unsupervised fashion. This problem shares a lot with the alignment likelihood in \textbf{Machine Translation}, even though here we can take advantage of the chunks sharing the same language (thus allowing some clever initialization of unsupervised methods with sensible guesses based, for instance, on Wordnet). On the other side our problem presents less precise training data, in that, while MT alignments are trained on sentence pairs, here we have a group of sentences sharing a same meaning; it can be argued that this cannot allow for strong syntactic assumptions.
% MT: T: more similar syntactic structure (label sentence by sentence)
\section{conclusions} \label{conclusions}
Here I described an algorithm to compare the similarity of two input sentences. Its main \textbf{advantages} are the concurrent consideration of many features like word meanings, position and overlapping, as well as the structure of the sentences, having the result in the form of aligned binary trees. The algorithm is recursive and exhaustive in that considers all the possible alignments of all the possible combinations of all the possible binary trees of the two input sentences, maximizing the efficiency with dynamic programming. Due to its modular structure, the algorithm can be easily extended with new features, allowing its integration with well-known Machine Learning techniques in Computational Linguistics (eg. Expectation-Maximization for computing word/phrase alignment probabilities).

There are some assumptions and \textbf{weak points} as well. First of all, the algorithm is inherently binary, in that every recursion step is made of a 2-split of the two input strings; this may have an effect in the way partial scores are combined. More generally, this "two times binary" recursive structure makes the features design task hard, since the effect of one operation is not easy to control over a potentially unlimited number of recursion steps. Also the number of parameters can become high, since each score is a linear combination of single features, requiring a separate Machine Learning training.

Lastly, this paper only describes the \textbf{backbone} of the algorithm, which is not yet ready for more formal tests on existing corpora. The main steps ahead before expecting good results from it include a better tuning and possible extension of the existing features  and the implementation of Machine Learning techniques to include alignment probabilities and more advanced statistical features for the chunks being compared (eg. their degree of informativeness).

%\begin{thebibliography}{9}

%\bibitem{achananuparp08}
%Palakorn Achananuparp, Xiaohua Hu, Xiajiong Shen. The Evaluation of Sentence Similarity Measures. Lecture Notes in Computer Science Volume 5182, 2008, pp 305-316

%\bibitem{udop}
%Rens Bod. Unsupervised parsing with U-DOP. CoNLL-X '06 Proceedings of the Tenth Conference on Computational Natural Language Learning, 2006, pp 85-92. 

%\bibitem{ibmmt}
%P. Brown, S. Della Pietra, V. Della Pietra, and R. Mercer (1993). The mathematics of statistical machine translation: parameter estimation. Computational Linguistics, 19(2), 263-311.

%\end{thebibliography}

%\input{./Chapters/Chapter4} 
%\input{./Chapters/Chapter5} 
%\input{./Chapters/Chapter6} 
%\input{./Chapters/Chapter7}

%----------------------------------------------------------------------------------------
%	THESIS CONTENT - APPENDICES
%----------------------------------------------------------------------------------------

\addtocontents{toc}{\vspace{2em}} % Add a gap in the Contents, for aesthetics

\appendix % Cue to tell LaTeX that the following 'chapters' are Appendices

% Include the appendices of the thesis as separate files from the Appendices folder
% Uncomment the lines as you write the Appendices

% Appendix A

\chapter{An Information State Update Example} % Main appendix title

\label{a:isu} % For referencing this appendix elsewhere, use \ref{isu}

\lhead{Appendix A. \emph{An Information State Update Example}} % This is for the header on each page - perhaps a shortened title

This appendix shows how SVPlay's Information State (IS) is updated during a \textbf{simple user interaction}, where the latter asks for the current song title. The interaction, at input/output level, is the following:
\begin{verbatim}
S> How can I help you?
U> what is the current song  
S> It's All Over Now, by The Rolling Stones
\end{verbatim}

As a premise, it is necessary to point out that, as OpenTDM comes with \textbf{no documentation}, the information contained in this appendix was derived from our analysis of the software's source code and log files. For this reason, it is possible that our explanation will not cover some parts of the process, either irrelevant or still obscure.

At t=0 the information state of the application is \textbf{empty}. While for the next steps we will only consider a subset of the OpenTDM's Total Information State (TIS), Figure \ref{a:isu:tis} shows the complete TIS, as it was reconstructed from the application's log. We can immediately spot, in the figure, the two main parts of the IS: private and shared. The most relevant parts of the \textbf{private} are
\begin{itemize}
	\item The list of the system's private \textbf{beliefs} (\texttt{bel})
	\item The stack of \textbf{plan} items that the system will execute.
\end{itemize}

From the \textbf{shared} part of the IS is important to highlight the stack of \textbf{questions under discussion} (\texttt{qud}), representing the issues that user and system will cooperatively try to solve. For the sake of record, we report that the \texttt{lu} field contains the \textit{last utterance} to be pronounced, either by the system or by the user; similarly, \texttt{pm} stands for \textit{previous move}, and \texttt{com} for \textit{common ground}.

\begin{figure}
\centering
$\left[\begin{tabular}{l l}
application &	\texttt{Application(`svplay', <Ontology>, <Domain>)} \\
device\_outputs &	\texttt{stack([])} \\
devices &	\texttt{\{`MplayDevice': <Device>\}} \\
input\_event &	\texttt{Event(START)} \\
latest\_moves & \texttt{open\_queue([`\#'])} \\
latest\_speaker & \texttt{None} \\
next\_utterance & 	$\left[\begin{tabular}{l l}
				alts & \texttt{[]} \\
				context & \texttt{\{\}} \\
				moves & \texttt{open\_queue([`\#'])} \\
				plan\_item & \texttt{None} \\
				\end{tabular}\right]$ \\
output & \  \\
passive\_mode & \texttt{False} \\
PRIVATE &	$\left[\begin{tabular}{l l}
			agenda & \texttt{open\_queue([`\#'])} \\
			bel & \texttt{\{\}} \\
			nim & \texttt{open\_queue([`\#'])} \\
			plan & \texttt{stack([do(top)])} \\
			\end{tabular}\right]$ \\
program\_state & \texttt{RUN} \\
recognised\_utterance & \texttt{None} \\
SHARED &	$\left[\begin{tabular}{l l}
			actions & \texttt{stackset([top])} \\
			com & \texttt{\{\}} \\
			issues & \texttt{stackset([])} \\
			lu &	$\left[\begin{tabular}{l l}
				moves & \texttt{\{\}} \\
				speaker & \texttt{None} \\
				turn\_count & \texttt{\{\}} \\
				\end{tabular}\right]$ \\
			pm & \texttt{\{\}} \\
			previous\_action & \texttt{None} \\
			qud & \texttt{stackset([])} \\
			\end{tabular}\right]$ \\
sys\_truns & 0 \\
timeout & 	$\left[\begin{tabular}{l l}
			duration & \texttt{None} \\
			enabled & \texttt{False} \\
			\end{tabular}\right]$ \\
\end{tabular}\right]$

\caption{An empty Total Information State in OpenTDM}
\label{a:isu:tis}
\end{figure}

We notice that \texttt{Event(START)} is pushed by default in the \texttt{input\_event} fileld (\texttt{application} and \texttt{device\_outputs} are filled as well, but they just hold static information, that will not change during the execution). This is because OpenTDM is \textbf{event-driven}: everything happening in the system is a consequence of an event, and \texttt{START} is the one that starts the causal chain. The complete list of events is defined in the \texttt{maharani.event} module.

As a consequence of the \texttt{START} event. the software will plan his first dialogue move. This is done by pushing the top plan in the \textbf{private plan} section. The top plan is defined in the application ontology;
\begin{verbatim}
["forget_all",
 "findout(?set([issue(?X.current_song(X)), issue(?X.my_name(X)),
	                action(increase_volume)]))"]
\end{verbatim}
This clears all the current believes and raises the issue of determine \textbf{what the user wants to do}. Literally, the \texttt{findout()} statement aims to decide whether the user wants to know the title of the current song, know the system's name or increase the volume. As a matter of fact, \pname supports more than these three actions, but the logic form of the top plan could not be changed accordingly because of the lack of documentation.

This plan has the effect of producing a \texttt{SYSTEM\_MOVES} event, that will prompt the user with the \textbf{question} ``How can I help you?". Note that the surface form of this question will be provided by the Language Unit, and is defined in \texttt{language/<APP><LANG>Prod.py}, in the application directory. Once the question is asked, it can be pushed in the ``questions under discussion" part of the \textbf{shared} IS.

The following is the resulting Information State:

\begin{table}[ht]
\small
$\left[\begin{tabular}{l l}
input\_event &	\texttt{Event(SYSTEM\_MOVES,} \\
\ & \texttt{[(Move(ask(?set([issue(?X.current\_song(X)),} \\
\ & \texttt{\hspace*{11em} issue(?X.my\_name(X)),} \\
\ & \texttt{\hspace*{11em} action(increase\_volume)])),} \\
\ & \texttt{\hspace*{3em} speaker=SYS, modality=speech),} \\
\ & \texttt{\hspace*{3em} `How can I help you?')])} \\
latest\_speaker & \texttt{SYS} \\
PRIVATE &	$\left[\begin{tabular}{l l}
			bel & \texttt{\{\}} \\
			plan & \texttt{stack([findout(?set([issue(?X.current\_song(X)),}\\
\ &			\texttt{\hspace*{5em}issue(?X.my\_name(X)),}\\
\ &			\texttt{\hspace*{5em} action(increase\_volume)]))])} \\
			\end{tabular}\right]$ \\
recognised\_utterance & \texttt{None} \\
SHARED &	$\left[\begin{tabular}{l l}
			com & \texttt{\{\}} \\
			qud & \texttt{stackset([?set([issue(?X.current\_song(X)),}\\
\ &			\texttt{\hspace*{5em}issue(?X.my\_name(X)),}\\
\ &			\texttt{\hspace*{5em} action(increase\_volume)]))])} \\
			\end{tabular}\right]$ \\
\end{tabular}\right]$
\end{table}

At this point, the \textbf{user} inputs the sentence ``What is the current song". This input is given to the language unit, that will return its formal interpretation, \texttt{ask(?X.current\_song(X))}. OpenTDM receives this interpretation (in the \texttt{tdm.interpret} module), and triggers a \texttt{USER\_MOVES} event.

This event causes the system to look for a \textbf{plan} in the domain that matches this user move. This plan exists, and its content is
\begin{verbatim}
dev_query(?X.current_song(X), MplayDevice)
\end{verbatim}
Which simply means to \textbf{query} the device for the title of the current song. Note that this new issue is also pushed in the shared questions under discussion.

The Information state is now the following:

\begin{table}[ht]
\small
$\left[\begin{tabular}{l l}
input\_event &	\texttt{Event(USER\_MOVES,} \\
\ & \texttt{open\_queue([Move(ask(?X.current\_song(X)),} \\
\ & \texttt{\hspace*{8em} speaker=USR, modality=speech), '\#']))} \\
latest\_speaker & \texttt{USR} \\
PRIVATE &	$\left[\begin{tabular}{l l}
			bel & \texttt{\{\}} \\
			plan & \texttt{stack([device\_query(?X.current\_song(X))])} \\
			\end{tabular}\right]$ \\
recognised\_utterance & \texttt{What is the current song} \\
SHARED &	$\left[\begin{tabular}{l l}
			com & \texttt{\{\}} \\
			qud & \texttt{stackset([?X.current\_song(X)]),} \\
\ &			\texttt{\hspace*{5em}?set([issue(?X.current\_song(X)),}\\
\ &			\texttt{\hspace*{8em}issue(?X.my\_name(X)),}\\
\ &			\texttt{\hspace*{8em}action(increase\_volume)]))}\\
			\end{tabular}\right]$ \\
\end{tabular}\right]$
\end{table}

The device answer to the query has the effect of creating a new \textbf{belief} in the private part of the IS:
\begin{verbatim}
current_song(`It's All Over Now, by The Rolling Stones')
\end{verbatim}
Which trivially states author and title of the song that the device is currently playing. This information is \textbf{output} with a \texttt{SYSTEM\_MOVES} event\footnote{Note that here the language unit receives the logical form \texttt{answer(current\_song(`It's All Over Now, by The Rolling Stones'))}, to be turned into a English sentence. In this case the logical form, that comes from the device, contains a complete English realization. A more elegant way of solving the problem would be to have the device output just the song name and author, and let the LU form the sentence.}.

After the system answer, the ``current song" issue is removed from the shared questions under discussion\footnote{Here the system also removes the top plan from \texttt{qud}; the reasons behind this choice are not clear.}, and the answer is added to the \textbf{common ground} (being the set of propositions that are considered to be in the beliefs of both the dialogue partners), along with the fact itself that the question has been answered.

The following is the information state after the system answer.

\begin{table}[ht]
\small
$\left[\begin{tabular}{l l}
input\_event &	\texttt{Event(SYSTEM\_MOVES, } \\
\ & \texttt{[(Move(answer(current\_song(`It's All Over Now, [...]')), } \\
\ & \texttt{\hspace*{8em} speaker=SYS, modality=speech), }\\
\ & \texttt{\hspace*{8em} `It's All Over Now, by The Rolling Stones')])} \\
latest\_speaker & \texttt{SYS} \\
PRIVATE &	$\left[\begin{tabular}{l l}
			bel & \texttt{current\_song(`It's All Over Now, by [...]')} \\
			plan & \texttt{stack([])} \\
			\end{tabular}\right]$ \\
recognised\_utterance & \texttt{What is the current song} \\
SHARED &	$\left[\begin{tabular}{l l}
			com & \texttt{\{current\_song(`It's All Over Now, by [...]'),}\\
\ &			\texttt{\hspace*{7em}resolved(?X.current\_song(X))\}} \\
			qud & \texttt{stackset([])}\\
			\end{tabular}\right]$ \\
\end{tabular}\right]$
\end{table}

After the felicitous utterance, the system will run out of plan elements to execute, and will thus load the \texttt{top plan} again.
%\input{./Appendices/AppendixB}
%\input{./Appendices/AppendixC}

\addtocontents{toc}{\vspace{2em}} % Add a gap in the Contents, for aesthetics

\backmatter

%----------------------------------------------------------------------------------------
%	BIBLIOGRAPHY
%----------------------------------------------------------------------------------------

\label{Bibliography}

\lhead{\emph{Bibliography}} % Change the page header to say "Bibliography"

\bibliographystyle{unsrtnat} % Use the "unsrtnat" BibTeX style for formatting the Bibliography

\bibliography{Bibliography} % The references (bibliography) information are stored in the file named "Bibliography.bib"

\end{document}  